\chapter{Methodology}
\label{chp:methodology}

% % \section{Notes}
% % In Brakerski's scheme, similar to previous FHE schemes based on LWE, the ciphertext $\overset{_\rightarrow}{c}$ and secret key $\overset{_\rightarrow}{s}$ are $n$-dimensional vectors whose dot product $\langle \overset{_\rightarrow}{c}, \overset{_\rightarrow}{s}\rangle \approx \mu$ equals the message $\mu$, up to some small ``error'' that is removed by rounding.

% \textcolor{red}{\textbf{Move implementation details from \autoref{sec:impl:hps} here?}}

\newcommand{\corelib}{\texttt{core} }
\newcommand{\serverlib}{\texttt{server} }
\newcommand{\clientlib}{\texttt{client} }

%%
%% Introduction
\vfill
A gap needed to be filled: implementing TFHE-style RGSW products (external and internal) in BGV-RNS. As a result of this thesis, a C++ library (\corelib) was developed to bring this functionality to the OpenFHE~library~\cite{cryptoeprint:2022/915}. The sPAR protocol has also been implemented in C++ using the RGSW library. I will stress the libraries are separate; the RGSW library can be used on its own as is, and the sPAR library can in theory be used without it --- provided a different library is developed to provide the same functionality and API. 

\begin{itemize}
    \item A result of this thesis: a C++ library which extends OpenFHE's BGV implementation to support RGSW operations
    \item And sPAR which was built on top of this library
    \item The code base is split into three sections found in the \texttt{lib} directory:
    \begin{enumerate}
        \item \corelib extends OpenFHE's \texttt{CryptoContext<T>} with the operations:
        \begin{itemize}
            \item \texttt{EncryptRGSW} outputs the RGSW encryption of the input plaintext
            \item \texttt{EvalExternalProduct} takes an RLWE and RGSW ciphertext
            \item \texttt{EvalInternalProduct} takes two RGSW ciphertexts
            \item Supporting functions like AddRGSW and MultRGSWPlaintext etc.
        \end{itemize}
        \item \serverlib builds the sPAR server functions on top of \corelib
        \item \clientlib builds the sPAR client functions on top of \corelib
    \end{enumerate}
    \item There are also two auxiliary directories \texttt{benchmarks} and \texttt{tests} that contain code to benchmark and unit test functions, respectively.
    \item And \texttt{scripts} contain python code for running the estimator
\end{itemize}

This section also explains the set up of the various tests; the numbers from running these tests are presented in \autoref{chp:results}.

\newpage

\paragraph{Bitwise Operations}
The library takes $b = \log_2 B$ as the parameter and derives $B = 2^b$, so that the arithmetic of digit extraction reduces to constant-time bitwise operations (\autoref{tab:bitwise}); in particular, reduction modulo $B$ becomes a mask and division by $B$ a shift.

\begin{table}[H]
    \caption{Bitwise Operators}
    \label{tab:bitwise}
    \centering
    \begin{tabular}{lcc}
        \toprule
        \textbf{Named Function} & \textbf{Operator Notation} & \textbf{Value}   \\
        \midrule
        \Call{LeftShift}{$x, b$}  & $x \ll b$     & $x\cdot 2^b$ \\
        \Call{RightShift}{$x, b$} & $x \gg b$     & $\lfloor x/2^b \rfloor$ \\
        \Call{And}{$x, y$}        & $x \land y$   & $\sum_i x_i^{(2)}\cdot y_i^{(2)} \cdot 2^i$ \\
        \Call{Modulo}{$x, b$}     & $x \land [(1 \ll b) - 1]$ & $x \bmod{2^b}$ \\
        \bottomrule
    \end{tabular}
\end{table}

%%%
%%% Center residue
%%%
\begin{algorithm}[ht!]
\caption[Center]{\protect\Call{Center}{$x,b$}}
\label{alg:center}
\begin{algorithmic}[1]
    \Require $x, b \in \mathbb{Z}$
    \Ensure $[x]_{2^b} \in [-2^{b-1}, 2^{b-1})$
    \State mask $\gets \Call{LeftShift}{1, b} - 1$
    \State half $\gets \Call{LeftShift}{1, b-1}$
    \State\Return $\Call{And}{x + \text{half}, \text{mask}} - \text{half}$
\end{algorithmic}
\end{algorithm}

%%
%% Implementation Details (RGSW)
%% The specifics of the program that are not necessarily important for the math, and other (novel) optimizations
%%
\newpage\section{Implementation Details of the Core Library}
\label{sec:impl:core}
The core library, located in the \texttt{libs/core} directory of the git repository \cite{gitrepo}, extends OpenFHE~\cite{cryptoeprint:2022/915} with \gls{rgsw} functionality that the framework does not itself expose. % Its centerpiece is the abstract interface \texttt{IExtendedContext}, which augments an OpenFHE \texttt{CryptoContext} with the essential \gls{rgsw} operations: encryption, the external product, and the internal product.
This functionality is specified by an abstract interface, \texttt{IExtendedContext}, that any project can consume independently of the underlying implementation; of which the essential operations are listed in \autoref{tab:main-functions}. The default implementation realizes the interface through gadget decomposition of the RNS primes, as described in \autoref{sec:gadget-decomposition}.

%% TODO: Call it BV-impl instead of default impl
% Either here or in gadget decomposition somewhere: Following the convention of~\cite{cryptoeprint:2021/204, cryptoeprint:2022/915}, we refer to this family of decomposition-based techniques collectively as \emph{BV}. 

\begin{table}[H]
    \centering
    \caption{Primary functions declared by \texttt{IExtendedContext}}
    \label{tab:main-functions}
    \begin{tabular}{lcc}
        \toprule
        \textbf{Function Signature}           & \textbf{Returns} \\ % & \textbf{Definition}  \\
        \midrule
        \Call{EncryptRGSW}{pk, $\mu$}         & $\Call{RGSW}{\mu}$ \\ % & \ref{def:rgsw} \\
        \Call{EvalExternalProduct}{$x$, $C$}  & $x \boxdot C$    \\ % & \ref{def:prod:external} \\
        \Call{EvalInternalProduct}{$A$, $C$}  & $A \boxtimes C$  \\ % & \ref{def:prod:internal} \\
        \bottomrule
    \end{tabular}
\end{table}

Following standard software engineering practice, the interface fixes \emph{what} each operation computes but not \emph{how}, so code written against it runs unchanged under any conforming implementation.

%% TODO: Explain context -> Takes crypto context parameters that are considered global/known by all the algorithms without specifying them!
%% TODO: Implementation is hard-coupled to the BGV-RNS scheme. Maybe: As OpenFHE provides a unified interface for both BGV and BFV --- and the TFHE products are defined independent of scheme --- it would not require much work to add support for BFV to the core library.

\paragraph{Bitwise Operations}
The library takes $b = \log_2 B$ as the parameter and derives $B = 2^b$, so that the arithmetic of digit extraction reduces to constant-time bitwise operations (\autoref{tab:bitwise}); in particular, reduction modulo $B$ becomes a mask and division by $B$ a shift.

\begin{table}[H]
    \caption{Bitwise Operators}
    \label{tab:bitwise}
    \centering
    \begin{tabular}{lcc}
        \toprule
        \textbf{Named Function} & \textbf{Operator Notation} & \textbf{Value}   \\
        \midrule
        \Call{LeftShift}{$x, b$}  & $x \ll b$     & $x\cdot 2^b$ \\
        \Call{RightShift}{$x, b$} & $x \gg b$     & $\lfloor x/2^b \rfloor$ \\
        \Call{And}{$x, y$}        & $x \land y$   & $\sum_i x_i^{(2)}\cdot y_i^{(2)} \cdot 2^i$ \\
        \Call{Modulo}{$x, b$}     & $x \land [(1 \ll b) - 1]$ & $x \bmod{2^b}$ \\
        \bottomrule
    \end{tabular}
\end{table}

%%%
%%% Center residue
%%%
\begin{algorithm}[ht!]
\caption[Center]{\protect\Call{Center}{$x,b$}}
\label{alg:center}
\begin{algorithmic}[1]
    \Require $x, b \in \mathbb{Z}$
    \Ensure $[x]_{2^b} \in [-2^{b-1}, 2^{b-1})$
    \State mask $\gets \Call{LeftShift}{1, b} - 1$
    \State half $\gets \Call{LeftShift}{1, b-1}$
    \State\Return $\Call{And}{x + \text{half}, \text{mask}} - \text{half}$
\end{algorithmic}
\end{algorithm}

\newpage\paragraph{Bitwise Operations}
The library takes $b = \log_2 B$ as the parameter and derives $B = 2^b$, so that the arithmetic of digit extraction reduces to constant-time bitwise operations (\autoref{tab:bitwise}); in particular, reduction modulo $B$ becomes a mask and division by $B$ a shift.

\begin{table}[H]
    \caption{Bitwise Operators}
    \label{tab:bitwise}
    \centering
    \begin{tabular}{lcc}
        \toprule
        \textbf{Named Function} & \textbf{Operator Notation} & \textbf{Value}   \\
        \midrule
        \Call{LeftShift}{$x, b$}  & $x \ll b$     & $x\cdot 2^b$ \\
        \Call{RightShift}{$x, b$} & $x \gg b$     & $\lfloor x/2^b \rfloor$ \\
        \Call{And}{$x, y$}        & $x \land y$   & $\sum_i x_i^{(2)}\cdot y_i^{(2)} \cdot 2^i$ \\
        \Call{Modulo}{$x, b$}     & $x \land [(1 \ll b) - 1]$ & $x \bmod{2^b}$ \\
        \bottomrule
    \end{tabular}
\end{table}

%%%
%%% Center residue
%%%
\begin{algorithm}[ht!]
\caption[Center]{\protect\Call{Center}{$x,b$}}
\label{alg:center}
\begin{algorithmic}[1]
    \Require $x, b \in \mathbb{Z}$
    \Ensure $[x]_{2^b} \in [-2^{b-1}, 2^{b-1})$
    \State mask $\gets \Call{LeftShift}{1, b} - 1$
    \State half $\gets \Call{LeftShift}{1, b-1}$
    \State\Return $\Call{And}{x + \text{half}, \text{mask}} - \text{half}$
\end{algorithmic}
\end{algorithm}

\newpage%% TODO: Consider including automatic B selection if there is time
% The base $B = \lfloor \max q_i \rfloor$ is chosen automatically so to be the smallest $B$ that covers the largest RNS prime with $\ell$ digits when the program starts. This makes the gadget exact, even if the $q_i$ are generated at different orders (in OpenFHE the first prime is typically larger than the other primes for some reason).
% Takes only $\ell$ as an input parameter and computes $\log{B} = \left\lfloor \max\{\log(q_i)\}/\ell \right\rfloor + 1$

\subsection{RGSW Encryption}
\begin{itemize}
    \item \texttt{ExtendedContextBV} implements the double decomposition; digit decomposing each RNS limb
    \item Let $P_g(v) = g\otimes v$ and $P_\mathsf{BV}(v) = P_g \circ \hps{P}(v) = g\otimes (\hps{P}(v))$
\end{itemize}
\begin{gather*}
    C = Z + \mu G = Z + \mu \cdot (I_2\otimes  P_g(\hps{P}(1)) = Z + (I_2\otimes  P_g(\hps{P}(\mu)) \\
    P_g(\hps{P}(\mu)) = g\otimes ([\mu]_{q_0}, \dots, [\mu]_{q_{k-1}}) \\
    ([\mu g]_{q_0}, \dots [\mu g]_{q_{k-1}}) \text{ where } \mu g = (\mu, \mu B, \dots, \mu B^{\ell})\\
    ([\mu]_{q_0}, [\mu B]_{q_0}, \dots, [\mu B^\ell]_{q_{k-1}}) \in \mathbb{Z}_Q^{k\cdot \ell}
\end{gather*}

\begin{itemize}
    \item Observe that this is the same information as $(\mu, \mu B, \dots, \mu B^{\ell})$
    \item Because the $i$-th tower of $\mu$ is $[\mu]_{q_i}$, the HPS gadget row $j$ gives $[[\mu]_{q_i}]_{q_j} = 0 \forall i\not = j$
\end{itemize}

\begin{gather*}
    %([\mu]_{q_0}, [\mu B]_{q_0}, \dots, [\mu B^\ell]_{q_{k-1}}) \\
    \left[ \begin{array}{ccc|ccc|ccc}
        [\mu]_{q_0} & \cdots & [\mu B^{\ell -1 }]_{q_0} & 0 & \dots & 0 & 0 & \dots & 0 \\
        0 &  \dots & 0 &[\mu]_{q_1} & \cdots & [\mu B^{\ell -1 }]_{q_{1}} & 0 & \dots & 0 \\
        \vdots & \ddots & \vdots & \vdots & \ddots & \vdots & \vdots & \ddots & \vdots \\
        0 & \dots & 0 & 0 & \dots & 0 &[\mu]_{q_{k-1}} & \cdots & [\mu B^{\ell -1 }]_{q_{k-1}}  \\
    \end{array} \right]
\end{gather*}

\begin{itemize}
    \item Ignore all the 0 elements, then we get the exact definition of the RNS form of $\mu$, $\mu B$, and so on
    \item However, each row of the RGSW must still encode only 1 tower from each of the $\ell$ polynomials
\end{itemize}

\begin{equation*}
    {P}_{\mathsf{BV}}(\mu) \;\overset{\tiny\text{flatten}}{\longmapsto}\; 
    \begin{bmatrix}
        [\mu]_{q_0} & \cdots & [\mu B^{\ell - 1}]_{q_0} \\
        [\mu]_{q_1} & \cdots & [\mu B^{\ell - 1}]_{q_1} \\
        \vdots & \ddots & \vdots  \\
        [\mu]_{q_{k-1}} & \cdots & [\mu B^{\ell - 1}]_{q_{k-1}}  \\
    \end{bmatrix}
    = ([\mu]_{\mathcal{B}_Q}, \dots, [\mu B^{\ell - 1}]_{\mathcal{B}_Q})
\end{equation*}

\begin{itemize}
    \item Multiplying $\mu$ by a factor of $B^i$ for a single $i$ can be done in NTT
    \item Applying this to $Z + P_\mathsf{BV}(\mu)$ allows parallelization on $k\ell$ threads
    \item $\mathcal{B}_Q$, $\ell$ and $B$ are all known/global variables tied to the context
\end{itemize}

%%%
%%% Encrypt RGSW-BV
%%%
\input{Algorithms/new/encrypt-bv}

\begin{itemize}
    \item Pre-computed value: $\Call{GetPower}{i, j} = B^i \bmod{q_j}$ where the modulo is just to keep numbers smaller in memory, negligible for performance 
    \item Then $\Call{TowerMul}{\mu, j, B^i}$ is really just extracting the $j-th$ limb, multiplying by the pre-computed value $B^i \bmod{q_j}$ calculate the product modulo $q_j$ before the tower is inserted back into $\mu$ 
\end{itemize}

\subsection{External Product}
%% TODO: Write about power-of-B gadget in earlier section!?
The decomposition on the power-of-$B$ gadget $g = (1, B, \dots, B^{\ell-1})$ is essentially digit extraction. The classical approach is to compute standard unsigned base-$B$ digits, producing a decomposition vector where each digit $d_i \in [0, B)$. Under our abstract gadget decomposition definition~\autoref{sec:gadget-decomposition}, the quality of this decomposition is $\beta = B - 1$.

From~\autoref{rem:extprod-asymmetry} we know the noise in the external product is a function of $\beta$. To exert tighter control over noise growth during external products, we naturally prefer to minimize this magnitude. By centering the digits into a signed representation where $d_i \in [-B/2, B/2)$, we halve the maximum absolute value, improving the overall decomposition quality to $\beta = B/2$. However, achieving this centered representation classically complicates parallelization. Standard signed recoding algorithms evaluate digits sequentially from least significant to most significant because shifting a positive digit $d_i \ge B/2$ into the negative range requires adding a $+1$ carry to the subsequent digit $d_{i+1}$. This ripple-carry dependency prevents us from slicing the integer into independent bit-windows, creating a computational bottleneck.

To avoid sequential carry propagation, we apply a well-known offset-binary trick by adding the constant $\varGamma$~\eqref{eq:carry-offset}.
\begin{equation}
    \label{eq:carry-offset}
    \varGamma = \frac{B}{2}\cdot\frac{B^{\ell} - 1}{B - 1}
\end{equation}

The second fraction is the standard mathematical identity for the sum of a finite geometric series, meaning~\eqref{eq:carry-offset} represents the $g$-decomposition of a base-$B$ number $\gamma$ where every digit is $B/2$.
\begin{gather*}
    \frac{B^{\ell} - 1}{B - 1} = \sum_{i=0}^{k-1} B^i = 1 + B + B^2 + \dots + B^{\ell - 1} \\
    \implies\;\; \varGamma = \sum_{i=0}^{k-1}\frac{B}{2}\cdot B^i = \langle \decomp{\gamma}{}, g \rangle
\end{gather*}

To decompose a positive integer $c \in \mathbb{Z}^{+}$ into signed digits $[d]_{B}$, we first apply the global offset to compute a shifted, non-negative integer $c' = c + \varGamma$. We then slice $c'$ into intermediate unsigned base-$B$ digits $u_i$. Because $B=2^b$ is a power of two, this extraction is performed purely via the independent bitwise \autoref{alg:bit:center}, completely bypassing sequential dependencies.

% \begin{algorithm}[ht!]
% \caption[Parallel Gadget Decomposition]{\protect\Call{ParallelDecompose}{$c, \ell, b$}}
% \label{alg:parallel-decomp}
% \begin{algorithmic}[1]
%     \Require Integer $c \in \mathbb{Z}$, decomposition length $\ell$, bit-length $b$
%     \Ensure Array of digits $d = (d_0, d_1, \dots, d_{\ell-1})$ where $d_i \in [-2^{b-1}, 2^{b-1})$ for $i < \ell - 1$
%     \State $\varGamma \gets \sum_{i=0}^{\ell-1} \Call{LeftShift}{1,\; b \cdot i + b - 1}$ \Comment{Precomputed constant offset}
%     \State $c' \gets c + \varGamma$ \Comment{Global shift}
%     \For{$i \gets 0$ \textbf{to} $\ell - 1$ \textbf{in parallel}}
%         \State $u_i \gets \Call{Modulo}{\Call{RightShift}{c',\; b \cdot i},\; b}$ \Comment{Extract independent bit-window}
%         \If{$i + 1 < \ell$}
%             \State $d_i \gets u_i - \Call{LeftShift}{1,\; b - 1}$ \Comment{Undo offset for lower digits}
%         \Else
%             \State $d_i \gets u_i$ \Comment{Top digit stays unsigned}
%         \EndIf
%     \EndFor
%     \State \Return $(d_0, d_1, \dots, d_{\ell-1})$
% \end{algorithmic}
% \end{algorithm}


% \begin{itemize}
%     \item \textbf{IMPORTANT:} decomposition determines the quality $\beta$, so the order of operations does matter for noise growth!
%     \item Decomposition used = per limb digit extraction
%     \item Must choose to decompose into $[-B/2, B/2)$ or $[0, B)$
%     \begin{itemize}
%         \item Centered has lower noise as the quality is $\beta = B/2$ but more complicated
%         \item Unsigned is far simpler and easily parallelizable but $\beta = B$
%     \end{itemize}
%     \item Centered representation creates an order; parallelization requires \textit{pre-processing}
% \end{itemize}


%%%
%%% Unsigned digit decomposition (optimized for B = 2^b)
%%% TODO: Move this to gadgets/BV section and have signed digit decomposition here
% \begin{algorithm}
% \caption{Decompose wrt. $B$ (\textcolor{red}{TODO: Signed digit decomposition})}
% \label{alg:bv:decompose}
% \begin{algorithmic}[1]
% \Require RNS polynomial $[x]_{\mathcal{B}_Q} \in \mathbb{Z}_{q_0} \times \mathbb{Z}_{q_1} \times \dots \times  \mathbb{Z}_{q_{k-1}}$, base $B = 2^b$, and decomposition parameter $\ell$
% \Ensure $\ell$ RNS ``digit'' polynomials $\mathbf{d} = (d_0, d_1, \dots, d_{k\ell-1}) \in \mathbb{Z}_B^{k\ell}$
% \State $a' \gets a$
% \For{$j = 0$ to $\ell - 1$}
%     \State $d_j \gets a' \pmod B$ \Comment{Extract the lower bits}
%     \If{$d_j \ge B/2$}
%         \State $d_j \gets d_j - B$ \Comment{Shift to signed range}
%     \EndIf
%     \State $a' \gets \frac{a' - d_j}{B}$ \Comment{Shift down and apply the carry}
% \EndFor
% \State \Return $\mathbf{d}$
% \end{algorithmic}
% \end{algorithm}

%%%
%%% 
%%%
% \begin{algorithm}
% \caption{Decompose wrt. $B$ (\textcolor{red}{TODO: Signed digit decomposition})}
% \label{alg:bv:decompose}
% \begin{algorithmic}[1]
% \Require Integer $[x]_{\mathcal{B}_Q} \in \mathbb{Z}_{q_0} \times \mathbb{Z}_{q_1} \times \dots \times  \mathbb{Z}_{q_{k-1}}$, base $B = 2^b$, and decomposition parameter $\ell$
% \Ensure $\ell$ RNS ``digit'' polynomials $\mathbf{d} = (d_0, d_1, \dots, d_{k\ell-1}) \in \mathbb{Z}_B^{k\ell}$
% \State $a' \gets a$
% \For{$j = 0$ to $\ell - 1$}
%     \State $d_j \gets a' \pmod B$ \Comment{Extract the lower bits}
%     \If{$d_j \ge B/2$}
%         \State $d_j \gets d_j - B$ \Comment{Shift to signed range}
%     \EndIf
%     \State $a' \gets \frac{a' - d_j}{B}$ \Comment{Shift down and apply the carry}
% \EndFor
% \State \Return $\mathbf{d}$
% \end{algorithmic}
% \end{algorithm}
\begin{algorithm}[ht!]
\caption[Parallel Signed Decomposition]{\protect\Call{Decompose}{$c, b, \ell$}}
\label{alg:decomp:parallel}
\begin{algorithmic}[1]
    \Require Integer $c \in \mathbb{Z}$, bit-length $b$ (where $B = 2^b$), decomposition parameter $\ell$
    \Ensure Array of signed digits $d = (d_0, d_1, \dots, d_{\ell-1})$ with quality $\beta = 2^{b-1}$
    
    \State $\varGamma \gets \sum_{i=0}^{\ell-1} 2^{b-1} \cdot 2^{b \cdot i}$ \Comment{This offset is typically precomputed}
    \State $c' \gets c + \varGamma$ \Comment{Apply the global shift}
    \State mask $\gets \Call{LeftShift}{1, b} - 1$
    \State half $\gets \Call{LeftShift}{1, b-1}$
    
    \For{$i \gets 0$ \textbf{to} $\ell - 1$ \textbf{in parallel}}
        \State $u_i \gets \Call{And}{\Call{RightShift}{c',\; i \cdot b},\; \text{mask}}$ \Comment{Extract the intermediate unsigned digit}
        \If{$i + 1 < \ell$}
            \State $d_i \gets u_i - \text{half}$ \Comment{Undo the offset to center the digit}
        \Else
            \State $d_i \gets u_i$ \Comment{The most significant digit remains unsigned}
        \EndIf
    \EndFor
    
    \State\Return $d = (d_0, d_1, \dots, d_{\ell-1})$
\end{algorithmic}
\end{algorithm}

Once we have a 

\begin{itemize}
    \item Uses the standard external product function \autoref{alg:external-product}
    \item Or a slightly parallelized version: the decompose function runs in parallel, not much gain from parallizing further, the additions are already cheap/fast-ish.
\end{itemize}

%%%
%%% External Product
%%%
\begin{algorithm}
\caption{External Product $\rlwe \boxdot \rgsw$ (Naive)}
\label{alg:external-product}
\begin{algorithmic}[1]
\Require RLWE ciphertext $(c_0, c_1) \in R_Q^2$, RGSW ciphertext $C \in R_Q^{2\ell}$
\Ensure RLWE ciphertext of product
\State $acc \gets (0, 0)$
\State $d_0 \gets \Call{Decompose}{c_0}$
\State $d_1 \gets \Call{Decompose}{c_1}$
\For{$r \gets [0, \ell)$}
    \State $acc[0] \gets acc[0] + d_0[r] \cdot C[r][0] + d_1[r]\cdot C[r+\ell][0]$
    \State $acc[1] \gets acc[1] + d_0[r] \cdot C[r][1] + d_1[r]\cdot C[r+\ell][1]$
\EndFor
\State \Return $acc$
\end{algorithmic}
\end{algorithm}

However, our implementation contains a series of optimizations aimed at . It can be divided into three distinct stages, each stage depending on the previous stage...

%%%
%%% Implemented (staged, parallel) external product
%%%
\input{Algorithms/new/extprod-bv}

\subsection{Internal Product}
The internal product is trivially implemented as repeated external products as per the definition~\autoref{def:prod:internal}.

%% TODO: Include algorithmic literal line-by-line as impl

\subsection{SIMD Batching}
Nothing in the gadget or in either product depends on how the plaintext ring is interpreted, so the implementation inherits the slot packing of \autoref{def:simd} from OpenFHE unchanged. A caller that builds its plaintexts with \texttt{MakePackedPlaintext} rather than \texttt{MakeCoefPackedPlaintext} obtains \gls{rgsw} ciphertexts whose message is a vector of $N$ independent slots, and every external and internal product then acts on all $N$ slots at once. The library selects the matching zero plaintext automatically, so no other change is required.

% The cost of a product is unchanged by this. A packed ciphertext is the same object of the same size, and the $2k\ell$ rows are decomposed and accumulated exactly as before, which is what makes the packing worthwhile. An application holding $N$ independent instances of the same binary computation can place one in each slot and evaluate all of them for the price of one, reducing the amortized cost per instance by a factor of $N$. The gain is in throughput and never in latency, and it is available only to applications whose work decomposes into independent instances in this way.
%%
%% Implementation Details (sPAR)
%% Specifics to the library, if any?
%%
\subsection{Implementation Details of sPAR}
\begin{itemize}
    \item Uses a procedural programming paradigme, meaning the library implements the required functions but does not provide a full implementation of the server nor client. 
    \item The benchmarks and testing framework is what actually composes the protocol, using the functions from this library
    \item This was done so that the functionality is not tightly coupled to the benchmarking without spending time implementing actual communication which would require additional work: serializing data, error handling, and coordinating multiple machines. 
    \item Does not implement the bucket sorting step: negligible impact
\end{itemize}

\subsection{Server Write Step}
\begin{itemize}
    \item Most important algorithm; does the actual homomorphic sorting
    \item Some changes compared to \cite{cryptoeprint:2025/860}
    \begin{itemize}
        \item Splits the write step and pre-processing step
        \item HasWrite is inverted i.e. ``notHasWrite''
        \item Uses external products to handle deep circuits as \cite{cryptoeprint:2025/860} recommends.
        \item Must use internal product
        \item \textbf{The server does not actually encrypt }
    \end{itemize}
\end{itemize}

%%%
%%% Server write loop
%%%
\begin{algorithm}
\caption{\texttt{Server::Write}}
\label{alg:spar:write}
\begin{algorithmic}[1]
\Require Reference to state matrices $I$ and $L$
\Require Encrypted RGSW value $V_r$
\Require RGSW indices $A_0$, $A_1$ and $A_2$ where $A_j$ is $\log{\eta}$ encrypted bits $\{c_{j,0}, \dots, c_{j,\eta -1}\}$ for $j \in [0, 3)$
\Ensure Updates the state matrices $I$ and $L$
\Ensure Encrypted boolean indicating wether the operation was successful
\State $\mathsf{hasNotWritten} \gets \Call{EncryptRGSW}{1}$
\For{$d \gets \{0, 1, 2\}$} % [0, D) 
    \For{$k \gets \{0, 1, 2\}$}
        \For{$i \gets [0, \eta)$}
            \State $h \gets z_{i,d}\boxtimes I_{i,k} \boxtimes \mathsf{hasNotWritten}$
            \State $L_{i,k} \gets L_{i,k} + V_r \boxdot h$
            \State $I_{i,k} \gets I_{i,k} - h$
            \State $\mathsf{hasNotWritten} \gets \mathsf{hasNotWritten} + h$
        \EndFor
    \EndFor
\EndFor
\State \Return $\mathsf{hasNotWritten}$
\end{algorithmic}
\end{algorithm}

\begin{remark}
    From a programming perspective, the choice of $K = 3$ and $D = 3$ is arbitrary. The code implements this as C++ generics also known as template specialization.  
\end{remark}
%
\newpage
\section{Experimental Setup}
\label{sec:experimental-setup}

\subsection{Protocol Simulation}
\begin{itemize}
    \item A single ``all-knowing'' orchestrator process (with access to all key material) simulates every party and performs the setup plus one full protocol round; no networking is involved.
    \item Round structure, parameterized over the number of users $\eta$ (\texttt{benchmark} directory):
    \begin{itemize}
        \item each virtual user holds a unique ID and a secret key share; the joint public key is produced by OpenFHE's chained multiparty key generation (each user extends the previous user's public key share; the final key is the joint public key);
        % FLAG: earlier notes said the public key shares are all "synthesized" into one key --- the generation is sequential/chained, not a combination of independent shares.
        \item each user submits an RLWE encryption of its value and $D$ one-hot RGSW index vectors under the joint key;
        % FLAG: earlier notes said each user "encodes their ID in an RGSW"; in the benchmark the written value is a plaintext value encrypted as RLWE, and the indices are the one-hot RGSW vectors of sec. impl:spar.
        \item the server executes \Call{Server::Write}{} once per user;
        \item decryption is multiparty: every user produces a partial decryption; fusion yields the plaintext.
    \end{itemize}
    \item The setup phase (key generation, state initialization) is considered free and excluded from all measurements.
\end{itemize}

\subsection{Unit Tests and Benchmarks}
\begin{itemize}
    \item Unit tests (\texttt{test} directory) empirically validate correctness: RGSW encryption and both products are tested for the plaintexts $0$, $1$ and a larger message, in both coefficient-packed and SIMD variants; the full protocol is tested parameterized over $\eta$ --- a passing run certifies the parameter set handles $\eta$ users.
    \item Tests check correctness at every stage and thereby localize failures (which operation failed and why), at the cost of speed.
    \item Benchmarks perform no correctness checks. One-shot timings suffice: client-side operations run independently per user, so the single-run time is the per-user latency, and for the server write only the per-user time is of interest.
    \item Tests and benchmarks share parameters and code paths, so a passing test validates exactly the computation the benchmark measures.
    \item The moduli $q_i$ are randomly generated per run; empirically the results do not depend on the specific $q_i$, only on the bit size of $Q$.
\end{itemize}

\subsection{Noise Measurement}
\begin{itemize}
    \item The testing environment always has access to the secret key, so RLWE noise is measured directly:
    \begin{equation*}
        x = (c_0, c_1) = (c_1 s + \Delta m + \epsilon,\; c_1)
        \implies \epsilon = c_0 - c_1 s - \Delta m, \qquad m = \Call{Decrypt}{x, s}.
    \end{equation*}
    % FLAG: state the reported metric explicitly (log2 of the infinity norm of eps, per limb?) --- ch. 4 tables depend on it.
    \item RGSW ciphertexts are intermediate objects; two ways to measure their noise:
    \begin{enumerate}
        \item extract the row encoding $\mu \cdot 1$ (exists whenever the gadget contains the entry $1$) and measure it as an RLWE ciphertext;
        \item compute the external product with a noiseless $\Call{RLWE}{1}$ (via \Call{MakePublicRGSW}{}-style zero-noise encryption) and measure the result.
    \end{enumerate}
    \item The latter is used: it applies to arbitrary gadgets, whereas the former assumes $1$ is a gadget entry --- traditional, but not guaranteed (and not true for the BV gadget rows scaled by $\tilde{e}_j$).
\end{itemize}
